\section{Introduction}\label{sec:intro}

\subsection{Background}
The colored HOMFLY-PT invariants of $\U(N)$ Chern--Simons theory are dual, at large $N$, to open topological strings on the resolved conifold, and the Ooguri--Vafa reformulation writes them in terms of integer BPS invariants (the LMOV invariants, after Labastida, Mari\~no, Ooguri and Vafa). For $\SO(N)$ and $\Sp(N)$ the dual theory is the topological string on an orientifold of the Calabi--Yau, as proposed in \cite{SV00}. It was conjectured in \cite{Marino09}, and verified there in many examples, that the colored Kauffman invariant of $\SO(N)$ minus one half of the $\U(N)$ invariant in composite representations is expressed through integer BPS invariants with crosscap number $c = 1, 2$ (see also \cite{BFM04,BFM05,NRZ13}). On the homological side, a Kauffman homology and superpolynomials for two-strand torus knots were proposed in \cite{GW05}, and the relation between $\SO$/$\Sp$ invariants, composite representations and BPS states was studied in \cite{GS11}.

Refined Chern--Simons theory was formulated in \cite{AgSh11}, and it was shown in \cite{KN17} that the $\U(N)$ refined Chern--Simons invariants produce refined BPS invariants on the resolved conifold which, for torus knots, are positive integers (refined large $N$ duality). In \cite{AgSc12} the refined topological string on an orientifold was defined as an M-theory index, the $\SO(2N)$ refined Chern--Simons invariants were computed, and the closed string free energy at large $N$ was shown to satisfy the integrality of second quantization. The closed $\SO$ and $\Sp$ free energies were organized with multiple sine functions in \cite{AM22}. For open strings, i.e.\ for knots, neither a formula extracting refined crosscap BPS invariants nor a test of their integrality has been available.

\subsection{Questions}
We address the following questions.
\begin{enumerate}
\item How does one extract a refined crosscap invariant from the $\SO(2N)$ refined Chern--Simons invariants?
\item Does it have an integral BPS expansion? If so, what structure does it have (universal factor, structure in the winding number $d$, boundary factors)?
\item Is positivity visible after refinement, i.e., do the coefficients align in sign in a suitable multiplet basis? Where does it hold and where does it break?
\item How is it related to the $\U(2N-1)$ cover, and how much of the crosscap structure is determined by the cover?
\item What remains open for higher representations?
\end{enumerate}
Question 1 is answered by the definition in \S\ref{sec:setup}. Question 2 is answered by the all-$d$ structure and the boundary factors (\S\ref{sec:index}) and by the integrality and the universal factor at general $N$ (\S\ref{sec:generalN}), including iterated torus knots. Question 3 is answered by the criterion in the fundamental representation (\S\ref{sec:torus}), by the signs of the brane--brane and brane--image sectors of the cover (\S\ref{sec:cover}) and by the multi-box crosscap signs (\S\ref{sec:positivity}): positivity holds in the fundamental representation of two-strand torus knots and on the cover side, and fails for the multi-box crosscap part and for the fundamental representation of knots that are not two-strand torus knots. Question 4 is addressed in \S\ref{sec:cover}, and question 5 in \S\ref{sec:positivity} and \S\ref{sec:discussion}.

\subsection{Main results}
Table~\ref{tab:summary} summarizes the range of verification.
\begin{itemize}
\item \textbf{Definition (\S\ref{sec:setup}).} A naive refinement of the bosonic kernel of \cite{Marino09} fails already for the unknot with two boxes: poles in $U = t^N$ do not cancel. With the dual Cauchy kernel of Koornwinder polynomials (the dual Cauchy identity of \cite{Mimachi01}) the product formula for the unknot is automatic. We take the uniform normalization in which the $\U(n)$ invariants are invariant under determinant shifts and charge conjugation, and the reformulation matrix $M(q^{-1})$ of \cite{Marino09}. The refined crosscap invariant is $g = f - \tfrac12 f^{\cov}$ (Definition \ref{def:crosscap}); $f$ is computed from the $\SO(2N)$ refined invariants and $f^{\cov}$ from the A-type invariants of the cover (ordinary and composite representations), separately and order by order in the number of boxes.
\item \textbf{All-$d$ structure (\S\ref{sec:index}).} The M-theory index of the orientifold gives $g = \tfrac12 E/(1+t) + C$ (Proposition \ref{prop:alld}), where $E$ is the signed sum over states invariant under the involution $\tilde I$ and $C$ the sum over frozen states; $E$ appears only in sectors with equal numbers of brane and image boxes. The same identity is the $\SO(2N)$ version of the trace representation of the Chern--Simons partition function of \cite{Witten11} (Remark \ref{rem:trace}). From the action of $\tilde I$ on boundary zero modes we derive the boundary factor $z_q^{n-1}(z_q^+)^n$ of $E$ (Proposition \ref{prop:boundary}) and verify it for the Hopf link and the framed unknot; for knots it is consistent with all computations. For closed strings the product formula of the $\SO(2N)$ partition function agrees with \cite{AgSc12}, and the ``Klein bottle'' term of \cite{AM22} is a rearrangement of this decomposition (Proposition \ref{prop:AM}). The separation of $E$ and $C$ is not determined by the index (Proposition \ref{prop:EC}).
\item \textbf{Unknot and Hopf link (\S\ref{sec:hopf}).} Integrality holds, the K\"ahler variable is shifted to $\tlam = \lambda(q/t)^{1/2}$, and $E_{22}$, $E_{24}$ are determined uniquely under a hypothesis on the numbers of states (Hypothesis \ref{hyp:QP}) and Hypothesis \ref{hyp:E2}.
\item \textbf{The criterion in the fundamental representation (\S\ref{sec:torus}).} For $T(2,2m+1)$, $m \le 5$, the $c = 1$ part has a $G$-expansion and the $c = 2$ part a closed form, with signs aligned in each $\lambda^\beta$ block. We state this as a criterion (Conjecture \ref{conj:criterion}). The criterion is specific to two-strand torus knots: for $T(3,4)$ and for algebraic cables the $c = 1$ part is not $\sigma$-symmetric, and for the cables the refined part of the $c = 2$ part has mixed signs (\S\ref{sec:cable}).
\item \textbf{Integrality and the universal factor at general $N$ (\S\ref{sec:generalN}).} Computing DAHA-Jones polynomials at generic $t$, we obtain the refined $\SO(2N)$ invariants at general $N$ for the trefoil with two to four boxes, $T(2,5)$ with two and three boxes, and $T(2,7)$, $T(2,9)$ with three boxes. The crosscap parts are integer Laurent polynomials in odd sectors, and $2(1+t)g_R$ is an integer Laurent polynomial in even sectors, the $(1+t)$ poles sitting exactly where Proposition \ref{prop:alld} predicts. The refined part is divisible by the universal factor $(\bt^2 - 1)(\ba^2 - \bq^2)$ (Conjecture \ref{conj:structure}).
\item \textbf{Algebraic cables (\S\ref{sec:cable}).} Taking the $D_N$ value of the iterated DAHA-Jones polynomial of \cite{CD14} as the definition, we obtain the fundamental representation and the two-box invariants of the $(2,13)$-cable of the trefoil and the fundamental representation of the $(2,15)$-cable at general $N$. Integrality, the position of the $(1+t)$ poles and the universal factor hold as for torus knots.
\item \textbf{Refined $\Sp(2N)$ (\S\ref{sec:sp}).} Implementing the $C_N$ DAHA, we obtain the refined $\Sp(2N)$ invariants of \cite{AM22} for the trefoil and $T(2,5)$ in the fundamental representation and for the trefoil with two boxes at general $N$. In the fundamental representation $g^{\Sp}(\ba) = -g^{\SO}(-\ba)$; with two boxes $W_R^{\Sp}(\ba) = (-1)^{|R|}W_R^{\SO}(-\ba)$, and the difference of the crosscap parts is the closed form $\pm\psi^2[\Delta]$. These are refined versions of the classical correspondence $\Sp(N) = \SO(-N)$ \cite{Cvitanovic08,MV10}.
\item \textbf{Structure of the cover (\S\ref{sec:cover}).} Splitting the cover into brane and image, the brane--brane sector coincides coefficient by coefficient with the $\U(N_A)$ refined LMOV invariants of \cite{KN17}. After removing boundary factors, the brane--image sector of two-strand torus knots is identified as a space of states built from $\Sym^k V$, $S \otimes S$ and universal corrections $F_{j'}$ (Result \ref{cr:principle}, Conjecture \ref{conj:closedform}); the generating function written with the rows of the fundamental and adjoint representations (Proposition \ref{prop:rows}) holds for fourteen torus knots. From this identification, fitted on $T(2,3)$--$T(2,9)$, $E$ is predicted for $T(2,11)$--$T(2,15)$ and reads as a signed partial sum of numbers of states.
\item \textbf{Multi-box status (\S\ref{sec:positivity}).} For higher representations the $c = 1$ part of the crosscap invariant is not $\sigma$-symmetric, has no $G$-expansion, and its signs do not align. The $\sigma$-asymmetry splits into a part that is already present in the unrefined theory, due to the parity rule relating $c$ to the power of $\lambda$ (in hatted invariants $\beta \equiv |R| - 1$), and a part specific to refinement, the $\sigma$-odd part divisible by $(\bt^2 - 1)(\ql - \ql^{-1})$. Since the cover is well behaved, the obstruction is localized in the multiplet structure of the frozen states.
\end{itemize}

\begin{table}[htbp]
\centering
\small
\caption{Range of verification. ``Integer'' means an integer Laurent polynomial; for even $|R|$ the crosscap part $g_R$ has the denominator $2(1+t)$ predicted by Proposition~\ref{prop:alld}, and ``$2(1+t)g_R$'' means that $2(1+t)g_R$ is integer while $g_{\ytwo} + g_{\yoneone}$ is integer as well. ``Universal factor'' is the divisibility of \S\ref{sec:universal}. ``Signs'': number of $\ba$-blocks of the quotient by the universal factor with aligned signs / total number of blocks; ``criterion'' refers to Conjecture~\ref{conj:criterion}; ``switching block'' is the block $k = 2m$ at which the universal series turns into the knot part (\S\ref{sec:split}).}\label{tab:summary}
\setlength{\tabcolsep}{4pt}
\begin{tabular}{p{4.3cm}p{2.7cm}p{2.5cm}p{5.0cm}}
\toprule
Object & Integer & Universal factor & Signs \\
\midrule
\addlinespace
Unknot, Hopf link ($a, b \le 4$) & yes & yes (10 components of $\hat C_{R,S}$) & $\hat g_{R,\yone}$ aligned; refined part of $\hat C_{R,S}$: 8/10 \\
\addlinespace
$T(2,2m+1)$, fundamental ($m \le 5$) & yes & yes & criterion holds \\
\addlinespace
$T(2,2m+1)$, fixed $N$ ($N = 2, 3$, $|R| \le 3$, $m \le 3$) & hatted invariants integer in all power-sum directions & --- & not decidable at fixed $N$ \\
\addlinespace
Trefoil, two boxes (large $N$) & $2(1+t)g_R$ & yes & 5/10 \\
\addlinespace
$T(2,5)$, two boxes & $2(1+t)g_R$ & yes & 6/18 \\
\addlinespace
Trefoil, three boxes & yes (no $(1+t)$ pole; hatted integer) & yes & 2/15, 3/14, 2/13 \\
\addlinespace
$T(2,5)$, $T(2,7)$, $T(2,9)$, three boxes & yes (no $(1+t)$ pole) & yes & 1/23--1/25, 1/33--1/35, 1/43--1/45 \\
\addlinespace
Trefoil, four boxes & $2(1+t)g_R$ ($(1+t)$ only in $p_4$, $p_2^2$) & yes & 0--2/18--22 \\
\addlinespace
$T(3,4)$, fundamental & yes & yes & criterion fails: $c = 1$ not $\sigma$-symmetric; $c = 2$ aligned per block; quotient of $g$: 4/4 \\
\addlinespace
Algebraic cables $\Cab{13}$, $\Cab{15}$, fundamental & yes & yes & criterion fails: $c = 1$ not $\sigma$-symmetric, refined part of $c = 2$ mixed; quotient of $g$: 6/9, 8/11 \\
\addlinespace
$\Cab{13}$, two boxes & $2(1+t)g_R$ & yes & 30/50 \\
\addlinespace
Cover, brane--brane (trefoil, $T(2,5)$, $|R| \le 2$) & yes (equals \cite{KN17}) & --- & all blocks aligned \\
\addlinespace
Cover, brane--image ($\Nfix$ with boundary factors removed; $T(2,3)$, $T(2,5)$, $T(2,7)$) & yes & --- & trefoil: all blocks aligned; $T(2,5)$, $T(2,7)$: aligned except the switching block \\
\addlinespace
Refined $\Sp(2N)$: $T(2,3)$, $T(2,5)$ fundamental; trefoil two boxes & yes (same structure as $\SO$) & --- & $g^{\Sp}(\ba) = -g^{\SO}(-\ba)$; $W_R$: $(-1)^{|R|}$ \\
\bottomrule
\end{tabular}
\end{table}

\subsection{Classification of claims and organization}
Statements are graded as follows. Propositions and corollaries are derived analytically, possibly under an explicitly stated hypothesis. A \emph{Result} is a statement verified by computer, with the range of verification stated. Conjectures are closed forms or structures suggested by computation beyond the verified range. Hypotheses are rules that could not be taken directly from the literature; we say which verification supports them. Definitions and remarks have their usual meaning. In the discussion, [derivation] and [computation] mark the two kinds of evidence, and \open\ marks an open point. Computational methods are described in Appendix \ref{app:code}. Abbreviations: DAHA (double affine Hecke algebra), LMOV (Labastida--Mari\~no--Ooguri--Vafa), TN (Taub--NUT), CoHA (cohomological Hall algebra).

\S\ref{sec:setup} fixes variables, kernel and normalization and defines the refined crosscap invariant. \S\ref{sec:index} derives the all-$d$ structure and the boundary factors from the orientifold index. \S\ref{sec:hopf} treats the unknot and the Hopf link, \S\ref{sec:torus} the criterion in the fundamental representation of torus knots, \S\ref{sec:generalN} the integrality and the universal factor at general $N$ for higher representations and for algebraic cables, and \S\ref{sec:sp} the refined $\Sp(2N)$ invariants. \S\ref{sec:cover} splits the cover into brane and image, \S\ref{sec:positivity} summarizes the status and obstructions for multi-box crosscap invariants, and \S\ref{sec:discussion} lists open problems.
